% Corrigé rendu compatible avec les recueils MiKTeX/T1.
\begin{corrigebox}[Corrigé extrait de la banque]

\CorrigeQuestion{1}
\begin{marginfigure}
\includegraphics[width=\linewidth]{14_01_C}
\end{marginfigure}

\CorrigeQuestion{2}
La fermeture géométrique s'écrit
\[
\vect{AB}+\vect{BC}+\vect{CA}=\vect{0},
\]
soit
\[
\lambda(t)\vect{i_2}-R\vect{i_1}-H\vect{j_0}=\vect{0}.
\]
Les projections dans $\rep{0}$ donnent
\[
\left\{
\begin{aligned}
\lambda(t)\cos\varphi(t)&=R\cos\theta(t),\\
\lambda(t)\sin\varphi(t)&=R\sin\theta(t)+H.
\end{aligned}
\right.
\]
D'où
\[
\tan\varphi(t)=
\frac{R\sin\theta(t)+H}{R\cos\theta(t)},
\]
et, en tenant compte du quadrant,
\[
\varphi(t)=\operatorname{atan2}
\bigl(R\sin\theta(t)+H,\,R\cos\theta(t)\bigr).
\]

\CorrigeQuestion{3}
La dérivation conduit à
\[
\boxed{
\dot{\varphi}(t)=
\frac{R\dot{\theta}(t)\left(R+H\sin\theta(t)\right)}
{R^2+H^2+2RH\sin\theta(t)}
}.
\]

\CorrigeQuestion{4}
Pour une fréquence de rotation de $10$ tours par minute,
\[
\omega=10\times\frac{2\pi}{60}=\frac{\pi}{3}\ \mathrm{rad\,s^{-1}}.
\]
En posant $\theta(t)=\omega t$, la fonction à tracer est
\[
\dot{\varphi}(t)=
\frac{R\omega\left(R+H\sin(\omega t)\right)}
{R^2+H^2+2RH\sin(\omega t)}.
\]

\begin{marginfigure}
\includegraphics[width=\linewidth]{14_02_c}
\end{marginfigure}

\end{corrigebox}
